\section{Why Convex-Hull Expansion}
\label{sec:hull}

The minimax objective \eqref{eq:minimax} inherits a sharp geometric structure in
the unrestricted limit $\mathcal{R} = \mathbb{R}^{d_g}$ (or equivalently, when the
restriction set is large enough that no face of $r$ binds at the optimum). In
this limit the partial-identification width is determined entirely by the
interaction between the null space of $\Sigma$ and the direction of the target
functional.

\subsection{Unrestricted dichotomy}
\label{sec:hull-dichotomy}

Write the linear target functional as $c(t)^\top \theta$ with
$c(t) := (\tilde c_m(t), \tilde c_r(t), 0, 0)$ (the zeros in the
$(\gamma_1, \gamma_2)$-coordinates follow because the target functional in
\eqref{eq:rep} is linear in $(\beta_1, \beta_2)$ only; the constant piece
$c_0(\gamma)$ shifts the endpoints but leaves the width invariant).

\begin{proposition}[Unrestricted width is $0$ or $\infty$]
\label{prop:dichotomy}
Let $\mathcal{R} = \mathbb{R}^{d_g}$. Then for any target direction $c(t)$ and
any feasible $b$,
\[
w_{\rm unr}(t; \Sigma, b) \;=\;
 \begin{cases}
  0 & \text{if } c(t)\in \mathrm{range}(\Sigma^\top), \\
  +\infty & \text{otherwise.}
 \end{cases}
\]
\end{proposition}

\begin{corollary}
\label{cor:collapse}
Adding a candidate site $k$ collapses the unrestricted width at target $t$ from
$\infty$ to $0$ if and only if the residualized context $(u_m^{(k)}, u_r^{(k)})$
of the candidate is such that
$c(t) \in \mathrm{range}\!\bigl(\Sigma_\mathcal{S}^\top + \Delta\Sigma_k^\top\bigr)$
but $c(t) \notin \mathrm{range}(\Sigma_\mathcal{S}^\top)$.
\end{corollary}

This is the unrestricted version of the \emph{coverage-expanding-donor} regime
(a) of the leave-one-site-out taxonomy in Boyarsky et al.\ (2026, \S4.4).

\subsection{The convex-hull interpretation}
\label{sec:hull-interp}

Proposition~\ref{prop:dichotomy} has a clean geometric reading in the site-level
context space. Let
\[
\mathcal{H}_\mathcal{S} \;:=\; \mathrm{conv}\!\bigl\{(C_m^{(s)}, C_r^{(s)})\,:\,
 s = 1, \ldots, S\bigr\} \;\subseteq\; \mathbb{R}^{d_M+d_I}.
\]
Its \emph{affine hull} $\mathrm{aff}(\mathcal{H}_\mathcal{S})$ is a linear
subspace (after centering) of dimension at most $S-1$. The cross-site moment
matrix $\Sigma_\mathcal{S}$ has rank at most $\dim\mathrm{aff}(\mathcal{H}_\mathcal{S})
 \cdot (\text{block factor})$, with equality when the existing sites have
sufficient overlap in $X$ at the points where their contexts differ.

\begin{proposition}[Hull expansion strictly reduces the null space]
\label{prop:hull}
If $(C_m^{(k)}, C_r^{(k)}) \notin \mathrm{aff}(\mathcal{H}_\mathcal{S})$ and the
candidate site's $X$-distribution has positive density on the individual-level
support, then
\[
\mathrm{rank}\bigl(\Sigma_{\mathcal{S}\cup\{k\}}\bigr)
 \;>\; \mathrm{rank}\bigl(\Sigma_\mathcal{S}\bigr),
\quad\Longleftrightarrow\quad
 \mathrm{null}\bigl(\Sigma_{\mathcal{S}\cup\{k\}}\bigr) \;\subsetneq\;
 \mathrm{null}\bigl(\Sigma_\mathcal{S}\bigr).
\]
\end{proposition}


Combined with Proposition~\ref{prop:dichotomy}, this yields a
finite-vs-infinite dichotomy: \emph{under no restriction $r$, the only candidate
sites that shrink the unrestricted width are those whose context lies outside the
affine hull of existing sites.} The convex-hull expansion rule is the special
case in which the target region $\mathcal{T}$ is itself contained in a larger
convex set whose existing-site coverage is incomplete.

\subsection{Width-as-distance-to-hull}
\label{sec:hull-distance}

Proposition~\ref{prop:dichotomy}'s dichotomy is discontinuous, which is
uninformative for comparing two candidates both of which lie outside the hull, or
both of which lie inside. The correct continuous refinement uses a
\emph{regularized} width,
\begin{equation}
\label{eq:regwidth}
w_r^\tau(t; \Sigma, b) \;:=\; \max_\theta\{c(t)^\top\theta\} -
 \min_\theta\{c(t)^\top\theta\}\quad
\text{s.t.}\quad
\Sigma\theta+b = 0,\; \|\theta\|_2 \le \tau,
\end{equation}
with a large-radius ball as a mild stand-in for a compact restriction set. The
regularized width admits a closed form:
\begin{equation}
\label{eq:distbound}
w_r^\tau(t; \Sigma, b) \;\le\; 2\tau \cdot \mathrm{dist}\!\bigl(
 c(t),\; \mathrm{range}(\Sigma^\top)\bigr),
\end{equation}
with equality up to a boundary adjustment that vanishes at large $\tau$. The
distance from $c(t)$ to $\mathrm{range}(\Sigma^\top)$ is precisely the norm of the
component of $c(t)$ in $\mathrm{null}(\Sigma)$. Geometrically, this distance is a
monotone function of the distance from the target profile $(t_m, t_r)$ to
$\mathrm{aff}(\mathcal{H}_\mathcal{S})$, and strictly monotone under the rank
conditions of Proposition~\ref{prop:hull}. We formalize this in
Proposition~\ref{prop:dist}.

\begin{proposition}[Width is monotone in distance to the affine hull]
\label{prop:dist}
Suppose the cross-site design is balanced and the within-site $X$-distribution is
the same across existing and candidate sites. Let
$d_\mathcal{S}(t) := \mathrm{dist}((t_m, t_r), \mathrm{aff}(\mathcal{H}_\mathcal{S}))$.
Then
\begin{equation}
\label{eq:prop3}
w_r^\tau(t; \Sigma_\mathcal{S}, b) \;=\; 2\tau \cdot d_\mathcal{S}(t) \cdot
 (1 + o_\tau(1)),
\end{equation}
with equality in the limit $\tau\to\infty$. In particular, for two candidates
$k_1, k_2$ with
$d_{\mathcal{S}\cup\{k_1\}}(t) < d_{\mathcal{S}\cup\{k_2\}}(t)$ uniformly over
$t\in\mathcal{T}$, candidate $k_1$ dominates $k_2$ under the minimax criterion
\eqref{eq:minimax} in the unrestricted limit.
\end{proposition}

Proposition~\ref{prop:dist} delivers the promised \textbf{convex-hull expansion
rule}: in the unrestricted (large-$\tau$) limit, the minimax-optimal candidate is
the one whose addition minimizes $\sup_{t \in \mathcal{T}}
d_{\mathcal{S}\cup\{k\}}(t)$ --- i.e., the one whose enlarged affine hull leaves
no target profile in $\mathcal{T}$ far from the span of sites. When $\mathcal{T}$
is a compact set containing points outside $\mathcal{H}_\mathcal{S}$, this
objective is exactly \emph{expanding the convex hull} in the direction of the
target region.
